\documentclass[11pt]{article}

\usepackage[margin=0.85in]{geometry}
\usepackage{amsmath}
\usepackage{graphicx}
\usepackage{hyperref}
\usepackage{enumitem}

\title{Constraint-Gated Waveform Design for Low-Excitation Ion Transport}
\author{Dan Hawkley}
\date{2026}

\begin{document}
\maketitle

\begin{abstract}
We present a compact computational pipeline for trapped-ion transport waveform design in segmented RF Paul traps. The pipeline maps electrode-basis potentials to moving potential wells, generates smooth voltage waveforms, simulates ion motion, and evaluates residual motional excitation. The initial implementation uses a transparent one-dimensional model so that the relationship between transport speed, waveform smoothness, and final excitation is directly visible.
\end{abstract}

\section{Introduction}

Segmented trapped-ion architectures require reliable ion transport between trap zones. Transport waveforms must move ions quickly while limiting residual motional excitation. This project develops a small, reproducible simulation pipeline for connecting waveform design to excitation metrics.

\section{Pipeline}

The computational pipeline is:
\[
\text{electrode basis}
\rightarrow
\text{potential well}
\rightarrow
\text{transport path}
\rightarrow
\text{voltage waveform}
\rightarrow
\text{ion motion}
\rightarrow
\text{excitation metric}.
\]

\section{Transport Model}

The first-pass model treats ion motion in a moving harmonic well:
\[
m\ddot{x}(t) = -m\omega^2[x(t)-x_c(t)],
\]
where \(x_c(t)\) is the waveform-controlled well center.

\section{Waveform Design}

Smooth transport paths are generated using a minimum-jerk interpolation:
\[
s(u)=10u^3-15u^4+6u^5.
\]
This produces zero velocity and acceleration at endpoints.

\section{Excitation Metric}

Residual excitation is estimated using final displacement and velocity relative to the target well:
\[
E_{\mathrm{res}} =
\frac{1}{2}m\left[v_f^2+\omega^2(x_f-x_c)^2\right].
\]

\section{Results}

Initial figures will compare target trajectories, ion trajectories, and residual excitation across transport durations.

\section{Discussion}

The project provides a bridge from trapped-ion transport literature into a reusable numerical workflow for waveform design, simulation, and constraint-gated evaluation.

\bibliographystyle{plain}
\bibliography{references}

\end{document}
